\chapter{Referencial Teórico}
\label{chap:theoretical-background}

Este capítulo apresenta os conceitos e tecnologias que
fundamentam o desenvolvimento da Medicano. São abordados
os temas de sistemas de informação em saúde, arquitetura
de software, tecnologias utilizadas, inteligência artificial
conversacional, computação em nuvem, qualidade de software
e metodologia de desenvolvimento.

\section{Sistemas de Informação em Saúde}
\label{sec:health-information-systems}

Os Sistemas de Informação em Saúde (SIS) são conjuntos de
componentes inter-relacionados que coletam, processam,
armazenam e comunicam dados para apoiar a tomada de decisão
e a gestão de processos no setor de saúde \cite{laudon2014}.
Na prática, esses sistemas abrangem desde prontuários
eletrônicos e sistemas de gestão hospitalar até plataformas
de agendamento e telemedicina.

No Brasil, a adoção de sistemas de informação em saúde
tem avançado de forma progressiva. O crescimento do
mercado de healthtechs — startups voltadas à digitalização
da saúde — reflete essa tendência: segundo
\citeonline{distrito2024}, o Brasil concentra 64,8\% das
healthtechs da América Latina, o que evidencia a
relevância e a maturidade do ecossistema nacional de
inovação em saúde.

A Medicano se insere nesse contexto como uma plataforma
que digitaliza o processo de agendamento de serviços de
saúde, integrando funcionalidades de busca, triagem e
confirmação de atendimentos em um único ambiente.

\section{Arquitetura de Software}
\label{sec:software-architecture}

\subsection{Arquitetura Cliente-Servidor}
\label{subsec:client-server}

A arquitetura cliente-servidor divide o sistema em dois
componentes com responsabilidades distintas: o cliente,
que lida com a interface e a interação com o usuário, e
o servidor, que processa as requisições, aplica as regras
de negócio e acessa o banco de dados \cite{sommerville2011}.

Essa separação traz vantagens importantes para projetos
como a Medicano. O frontend e o backend podem ser
desenvolvidos, implantados e escalados de forma
independente, o que facilita a manutenção e permite que
cada camada evolua no seu próprio ritmo sem impactar
a outra \cite{sommerville2011}.

\subsection{Arquitetura REST}
\label{subsec:rest}

O Representational State Transfer (REST) é um estilo
arquitetural para sistemas distribuídos definido por
\citeonline{fielding2000} em sua dissertação de doutorado.
Ele estabelece um conjunto de restrições para a comunicação
entre cliente e servidor, das quais as mais relevantes
para aplicações web são a ausência de estado nas
requisições (\textit{stateless}) e a interface uniforme
baseada nos métodos HTTP.

Em uma API RESTful, cada recurso é identificado por uma
URI e as operações sobre ele são realizadas por meio dos
métodos HTTP padrão: GET para leitura, POST para criação,
PUT e PATCH para atualização e DELETE para remoção
\cite{fielding2000}. Essa padronização facilita o
consumo da API por qualquer cliente, independentemente
da tecnologia utilizada.

A Medicano adota o estilo REST em toda a sua API,
permitindo que o frontend React consuma os serviços
do backend NestJS de forma padronizada e desacoplada.

\subsection{Padrão MVC}
\label{subsec:mvc}

O padrão Model-View-Controller (MVC) organiza o código
de uma aplicação em três camadas com responsabilidades
bem definidas: o Model, que representa os dados e as
regras de negócio; a View, que cuida da apresentação
para o usuário; e o Controller, que recebe as entradas,
aciona o Model e decide qual View exibir
\cite{sommerville2011}.

Essa separação evita que lógica de negócio fique misturada
com código de interface, tornando o sistema mais fácil
de manter e testar \cite{martin2017}. No backend da
Medicano, o NestJS implementa uma variação desse padrão
por meio de controllers, que recebem as requisições HTTP,
e services, que concentram toda a lógica de negócio da
aplicação.

\section{Tecnologias Utilizadas}
\label{sec:technologies}

\subsection{Node.js}
\label{subsec:nodejs}

O Node.js é um ambiente de execução JavaScript no lado
do servidor construído sobre o motor V8 do Google Chrome.
Ele utiliza um modelo de I/O não bloqueante orientado a
eventos, o que o torna eficiente para aplicações que
precisam lidar com muitas conexões simultâneas sem
consumir grandes quantidades de memória
\cite{tilkov2010}.

Na prática, isso significa que o Node.js processa
requisições de forma assíncrona em uma única thread,
sem bloquear a execução enquanto aguarda respostas de
banco de dados ou serviços externos. Essa característica
é especialmente relevante para a Medicano, cuja API
precisa lidar com requisições de agendamento e chamadas
a modelos de linguagem ao mesmo tempo \cite{tilkov2010}.

\subsection{NestJS}
\label{subsec:nestjs}

O NestJS é um framework para construção de APIs em
Node.js que usa TypeScript como linguagem principal.
Ele organiza o código em módulos independentes, onde
cada módulo representa um domínio da aplicação — como
autenticação, clínicas ou agendamentos — com seus
próprios controllers, services e providers
\cite{nestjs2024}.

Essa organização modular torna o projeto mais fácil
de manter e testar, já que cada parte pode ser
desenvolvida e modificada de forma isolada. Além disso,
o NestJS já vem com suporte nativo a validação de dados,
documentação de API com Swagger e testes com Jest, o
que reduz a quantidade de configuração necessária para
atingir um nível profissional de qualidade de software
\cite{nestjs2024}.

Na Medicano, o NestJS foi escolhido em substituição ao
Express justamente por oferecer essa estrutura pronta.
Enquanto o Express exige que o desenvolvedor defina
a organização do projeto do zero, o NestJS impõe boas
práticas arquiteturais desde o início, o que é
especialmente importante em um projeto desenvolvido
por uma equipe.

\subsection{React}
\label{subsec:react}

O React é uma biblioteca JavaScript para construção de
interfaces de usuário desenvolvida e mantida pela Meta.
Ele é baseado no conceito de componentes reutilizáveis
e permite a construção de Single Page Applications
(SPAs), nas quais a interface é atualizada dinamicamente
sem a necessidade de recarregar a página inteira
\cite{reactdev2024}.

O React utiliza um DOM virtual — uma representação em
memória da interface — para identificar apenas o que
mudou entre um estado e outro e aplicar somente essas
alterações na tela. Isso resulta em interfaces mais
rápidas e responsivas \cite{reactdev2024}. Na Medicano,
o React é utilizado em conjunto com TypeScript e Vite
para construção de todas as telas da plataforma.

\subsection{TypeScript}
\label{subsec:typescript}

O TypeScript é um superset tipado do JavaScript
desenvolvido pela Microsoft que adiciona tipagem
estática opcional à linguagem \cite{microsoft2012}.
Na prática, isso significa que erros como acessar uma
propriedade inexistente em um objeto ou passar um
argumento do tipo errado para uma função são detectados
em tempo de compilação, antes mesmo de o código ser
executado.

Pesquisas indicam que a adoção de tipagem estática
reduz em até 15\% os bugs que chegam a produção em
projetos JavaScript \cite{gao2017}. Além disso, o
TypeScript melhora significativamente a experiência
de desenvolvimento com recursos como autocompletar
e refatoração segura nas IDEs. Na Medicano, o TypeScript
é utilizado tanto no frontend quanto no backend,
garantindo consistência de tipos em toda a aplicação.

\subsection{Bancos de Dados NoSQL e MongoDB}
\label{subsec:mongodb}

Os bancos de dados NoSQL surgiram como alternativa aos
bancos relacionais para situações em que a flexibilidade
de schema e a escalabilidade horizontal são mais
importantes do que a consistência transacional rígida
\cite{sadalage2012}. Em vez de tabelas com colunas
fixas, os bancos NoSQL podem armazenar dados em
formatos como documentos JSON, grafos ou pares
chave-valor.

O MongoDB é um banco de dados NoSQL orientado a
documentos que armazena dados em formato BSON (Binary
JSON) \cite{sadalage2012}. Essa flexibilidade é
especialmente útil para a Medicano, cujo modelo de
dados inclui entidades com estruturas variáveis —
como os slots de disponibilidade dos profissionais,
que diferem entre especialidades e tipos de
atendimento. O Mongoose é utilizado como ODM (Object
Document Mapper) para definir schemas, validações e
relacionamentos entre os documentos de forma
declarativa no NestJS.

\subsection{JSON Web Token}
\label{subsec:jwt}

O JSON Web Token (JWT) é um padrão aberto definido pela
RFC 7519 para transmissão segura de informações entre
partes como um objeto JSON \cite{rfc7519}. Um token JWT
é formado por três partes separadas por pontos: o
cabeçalho, que define o algoritmo de assinatura; o
payload, que contém os dados do usuário; e a assinatura,
que garante que o token não foi alterado.

Por ser stateless, o servidor não precisa armazenar
sessões — a cada requisição o token é enviado pelo
cliente e validado pela assinatura digital, sem
necessidade de consultar o banco de dados \cite{rfc7519}.
Isso torna o JWT adequado para APIs REST consumidas por
frontends desacoplados, como é o caso da Medicano, onde
o React se comunica com o NestJS em domínios distintos.

\subsection{Redis e Gerenciamento de Sessões}
\label{subsec:redis}

O Redis é um armazenamento de dados em memória de
código aberto que pode ser utilizado como banco de
dados, cache e broker de mensagens \cite{redis2024}.
Por operar inteiramente na memória RAM, o Redis oferece
latência extremamente baixa para operações de leitura
e escrita, tornando-o adequado para casos de uso que
exigem acesso frequente e rápido a dados temporários.

Uma das funcionalidades mais relevantes do Redis para
aplicações web é o suporte nativo a TTL (Time to Live),
que permite definir um tempo de expiração automática
para cada chave armazenada \cite{redis2024}. Na
Medicano, o Redis é utilizado para o gerenciamento
do ciclo de vida dos tokens JWT. Após o login, o token
gerado é armazenado no Redis com um TTL correspondente
ao tempo de expiração do token. A cada requisição
autenticada, além da validação criptográfica da
assinatura JWT, o sistema verifica se o token ainda
existe no Redis. No logout, o token é imediatamente
removido do Redis, invalidando-o antes do seu prazo
de expiração natural. Essa abordagem resolve a
limitação inerente ao JWT stateless, que por padrão
não permite a revogação de tokens antes do vencimento
\cite{rfc7519}.

\subsection{Bcrypt}
\label{subsec:bcrypt}

O Bcrypt é um algoritmo de hashing projetado
especificamente para armazenamento seguro de senhas
\cite{provos1999}. Diferentemente de funções de hash
genéricas como MD5 ou SHA-1, o Bcrypt é intencionalmente
lento — ele possui um fator de custo configurável que
determina quantas iterações o algoritmo executa, tornando
ataques de força bruta computacionalmente inviáveis mesmo
com hardware moderno.

O Bcrypt também incorpora um salt aleatório em cada hash,
o que impede o uso de tabelas pré-computadas para reverter
os hashes \cite{provos1999}. Na Medicano, o Bcrypt é
utilizado para armazenar as senhas dos usuários de forma
que, mesmo em caso de vazamento do banco de dados, as
senhas originais não possam ser recuperadas.

\section{Inteligência Artificial Conversacional}
\label{sec:conversational-ai}

\subsection{Large Language Models}
\label{subsec:llm}

Os Large Language Models (LLMs) são modelos de linguagem
baseados em redes neurais profundas treinados em grandes
volumes de texto para aprender padrões da linguagem humana.
Segundo \citeonline{brown2020}, modelos dessa categoria
são capazes de realizar tarefas de processamento de
linguagem natural sem treinamento específico para cada
tarefa, utilizando apenas exemplos fornecidos no contexto
da requisição — técnica conhecida como
\textit{few-shot learning}.

Na prática, isso significa que um LLM pode conduzir uma
conversa, responder perguntas, resumir textos e seguir
instruções complexas sem precisar ser retreinado para
cada novo caso de uso. Essa versatilidade os torna
adequados para aplicações conversacionais como o chatbot
de triagem da Medicano, onde o modelo precisa interpretar
sintomas descritos em linguagem livre e recomendar uma
especialidade de saúde \cite{brown2020}.

\subsection{Engenharia de Prompt}
\label{subsec:prompt-engineering}

A engenharia de prompt é a prática de projetar as
instruções fornecidas a um LLM para obter respostas mais
precisas e alinhadas ao objetivo da aplicação
\cite{wei2022}. Em vez de treinar ou ajustar o modelo,
o desenvolvedor define um \textit{system prompt} que
estabelece o comportamento esperado, o tom das respostas
e as restrições do sistema.

Em aplicações de saúde, a qualidade do prompt é
especialmente crítica. \citeonline{wei2022} demonstram
que técnicas como \textit{chain-of-thought} — instruir
o modelo a raciocinar passo a passo antes de responder
— aumentam significativamente a precisão em tarefas
que exigem raciocínio estruturado.

No chatbot de triagem da Medicano, o system prompt
instrui o modelo a conduzir uma conversa guiada sobre
os sintomas do usuário e, ao final, recomendar uma
das especialidades disponíveis na plataforma. O prompt
também inclui restrições explícitas que impedem o
modelo de emitir diagnósticos médicos ou prescrever
medicamentos, limitando sua atuação ao direcionamento
para o profissional adequado.

\subsection{Vercel AI SDK}
\label{subsec:vercel-ai-sdk}

O Vercel AI SDK é uma biblioteca TypeScript de código
aberto para integração de LLMs em aplicações web e APIs
\cite{vercelaisdk2024}. Ele oferece uma interface
unificada e independente de provedor, permitindo que a
aplicação utilize modelos da OpenAI, Anthropic, Google
e outros com mínima alteração de código.

Entre os principais recursos do SDK estão o suporte
nativo a streaming de respostas — que permite exibir
o texto gerado progressivamente enquanto ainda está
sendo processado — e o gerenciamento do histórico de
conversa, que mantém o contexto entre as mensagens
\cite{vercelaisdk2024}. Na Medicano, o Vercel AI SDK
é integrado ao módulo de chat do NestJS e é responsável
pela comunicação com o modelo de linguagem e pelo
streaming das respostas para o frontend React.

A escolha do Vercel AI SDK em detrimento de alternativas
como o LangChain se justifica pela menor complexidade
de configuração e pelo suporte nativo a TypeScript, que
se alinha com a stack do projeto. O LangChain, embora
mais poderoso para casos como agentes com múltiplas
ferramentas e RAG, adiciona complexidade desnecessária
para o escopo do MVP \cite{vercelaisdk2024}.

\subsection{Considerações Éticas e Regulatórias}
\label{subsec:ai-ethics}

A utilização de inteligência artificial em aplicações
de saúde está sujeita a considerações éticas e
regulatórias importantes. O Conselho Federal de Medicina,
por meio da Resolução CFM n.º 2.336/2023, estabelece
diretrizes para o uso de tecnologias digitais no
atendimento médico, incluindo restrições sobre
diagnóstico automatizado \cite{cfm2023}.

O chatbot de triagem da Medicano é projetado
explicitamente como uma ferramenta de orientação e
direcionamento, e não de diagnóstico. Toda interação
inclui um aviso claro ao usuário de que as recomendações
geradas não substituem a avaliação de um profissional
de saúde habilitado. Além disso, os dados compartilhados
nas conversas são classificados como dados sensíveis
pela Lei Geral de Proteção de Dados (LGPD) e tratados
com as salvaguardas correspondentes, incluindo
criptografia em trânsito e restrição de acesso
\cite{cfm2023}.

\section{Computação em Nuvem}
\label{sec:cloud}

\subsection{Infraestrutura como Serviço}
\label{subsec:iaas}

A computação em nuvem é definida por
\citeonline{vaquero2009} como um conjunto de recursos
virtualizados facilmente acessíveis que podem ser
dinamicamente reconfigurados para atender a cargas
variáveis. O modelo de Infraestrutura como Serviço
(IaaS) é uma das camadas fundamentais da computação
em nuvem, na qual o provedor oferece servidores
virtuais, armazenamento e rede como serviços sob
demanda, sem que o cliente precise gerenciar o
hardware físico subjacente \cite{vaquero2009}.

O Amazon EC2 é o serviço de IaaS da Amazon Web Services
que permite provisionar servidores virtuais com
configurações flexíveis de CPU, memória e armazenamento
\cite{awsec2}. Na Medicano, uma instância EC2 hospeda
o backend NestJS, oferecendo controle total sobre o
ambiente de execução e permitindo a configuração do
Nginx como proxy reverso.

\subsection{Armazenamento de Objetos e CDN}
\label{subsec:s3-cloudfront}

O Amazon S3 é um serviço de armazenamento de objetos
da AWS que, quando configurado para hospedagem de
sites estáticos, permite servir arquivos HTML, CSS e
JavaScript diretamente para os usuários sem a
necessidade de um servidor web dedicado \cite{awss3}.

O Amazon CloudFront é uma rede de distribuição de
conteúdo (CDN) que entrega o conteúdo do S3 a partir
de pontos de presença geograficamente distribuídos,
reduzindo a latência para usuários em diferentes
regiões \cite{awscf}. Na Medicano, o build estático
do frontend React é hospedado no S3 e distribuído
via CloudFront, o que resulta em menor latência,
maior disponibilidade e custo operacional reduzido
em comparação com a hospedagem em uma segunda
instância EC2.

\subsection{Proxy Reverso com Nginx}
\label{subsec:nginx}

O Nginx é um servidor web de alto desempenho que
também pode ser utilizado como proxy reverso,
recebendo as requisições dos clientes e encaminhando
para o servidor de aplicação interno \cite{reese2008}.
Essa configuração adiciona uma camada de segurança
entre a internet e a aplicação, impedindo que o
servidor NestJS fique diretamente exposto.

Na Medicano, o Nginx é configurado na instância EC2
para receber requisições na porta 443 com HTTPS,
realizar a terminação SSL e encaminhar as requisições
para o NestJS na porta 3000. Os certificados SSL são
gerados e renovados automaticamente pelo Let's Encrypt,
garantindo a criptografia do tráfego sem custo adicional
\cite{reese2008}.

\subsection{Gerenciamento de Segredos}
\label{subsec:secrets}

O armazenamento de credenciais diretamente no
código-fonte ou em arquivos de configuração versionados
representa um risco significativo de segurança, uma
vez que qualquer pessoa com acesso ao repositório
passa a ter acesso às credenciais de produção
\cite{awssm}.

O AWS Secrets Manager é um serviço gerenciado da AWS
para armazenamento e recuperação segura de credenciais
\cite{awssm}. Na Medicano, todas as informações
sensíveis — incluindo a string de conexão do MongoDB,
o segredo JWT e as chaves de API dos modelos de
linguagem — são armazenadas no Secrets Manager e
recuperadas pela aplicação em tempo de execução por
meio do SDK da AWS, seguindo o princípio do menor
privilégio.

\section{Qualidade de Software}
\label{sec:software-quality}

\subsection{Validação de Dados e DTOs}
\label{subsec:validation}

A validação de dados de entrada é um aspecto fundamental
da segurança e da integridade de qualquer aplicação web.
O padrão Data Transfer Object (DTO) define a estrutura
dos dados que trafegam entre as camadas da aplicação,
servindo como contrato entre o cliente e o servidor
\cite{martin2017}.

No NestJS, os DTOs são implementados como classes
TypeScript decoradas com anotações do pacote
\texttt{class-validator}, que permitem definir regras
de validação diretamente na classe \cite{nestjs2024}.
Quando o pipe de validação global está ativo, qualquer
requisição com dados fora do contrato é automaticamente
rejeitada com uma resposta HTTP 400 antes de chegar
à camada de serviço, protegendo a aplicação contra
dados malformados.

\subsection{Contêinerização com Docker}
\label{subsec:docker}

O Docker é uma plataforma de contêinerização que
permite empacotar uma aplicação e suas dependências
em um contêiner isolado e reproduzível
\cite{docker2024}. No contexto do desenvolvimento
da Medicano, o Docker é utilizado exclusivamente
para padronizar o ambiente de infraestrutura local
— MongoDB e Redis — garantindo que todos os membros
da equipe trabalhem com as mesmas versões e
configurações, independentemente do sistema
operacional da máquina.

O Docker Compose orquestra a subida dos dois
serviços com um único comando, simplificando a
configuração do ambiente e eliminando a necessidade
de instalação manual de dependências de
infraestrutura nas máquinas de desenvolvimento.

\subsection{Testes Unitários}
\label{subsec:unit-tests}

Os testes unitários verificam o comportamento de
unidades individuais de código de forma isolada,
garantindo que cada componente funcione corretamente
de acordo com sua especificação \cite{sommerville2011}.
A prática de testes unitários reduz o custo de
correção de bugs, pois os defeitos são encontrados
antes de chegarem a produção.

O Jest é um framework de testes JavaScript desenvolvido
pela Meta com suporte nativo a mocks e relatórios de
cobertura de código \cite{jestjs2024}. O NestJS inclui
suporte nativo ao Jest em sua configuração padrão,
facilitando a escrita de testes para controllers e
services. Na Medicano, os testes unitários são
aplicados principalmente à camada de services, onde
está concentrada a lógica de negócio da aplicação.

\subsection{Documentação de API com OpenAPI}
\label{subsec:openapi}

A especificação OpenAPI é um padrão aberto para
descrição de APIs RESTful que permite gerar
documentação interativa automaticamente a partir
das definições da API \cite{swagger2021}. Essa
documentação facilita o consumo da API por outras
equipes e serve como referência durante o
desenvolvimento do frontend.

O NestJS oferece integração nativa com o Swagger
por meio do pacote \texttt{@nestjs/swagger}, que
utiliza decorators TypeScript para extrair os
metadados dos controllers e DTOs e gerar a
documentação automaticamente \cite{nestjs2024}.
Na Medicano, a documentação Swagger é acessível
pelo endpoint \texttt{/api/docs} e serve como
base para a especificação completa da API descrita
no Apêndice~\ref{appendix:api-spec}.

\section{Metodologia de Desenvolvimento}
\label{sec:dev-methodology}

\subsection{Desenvolvimento Ágil}
\label{subsec:agile}

O desenvolvimento ágil é um conjunto de práticas
baseadas nos valores do Manifesto Ágil, publicado
em 2001, que prioriza a entrega contínua de software
funcionando, a colaboração com o cliente e a
capacidade de responder a mudanças ao longo do
projeto \cite{agilemanifesto2001}.

O Scrum é o framework ágil mais adotado na indústria
de software \cite{schwaber2020}. Ele organiza o
trabalho em ciclos curtos denominados sprints, com
duração de uma a quatro semanas, ao final dos quais
um incremento funcional do produto deve estar
disponível. Segundo \citeonline{schwaber2020}, o
Scrum se baseia em três pilares: transparência,
inspeção e adaptação — o que permite que a equipe
identifique e corrija desvios rapidamente.

Na Medicano, o desenvolvimento é organizado em
sprints quinzenais, com revisões ao final de cada
ciclo para avaliar o progresso e ajustar as
prioridades do backlog conforme necessário.

\subsection{Conceito de MVP}
\label{subsec:mvp}

O Produto Mínimo Viável (MVP) é definido por
\citeonline{ries2011} como a versão de um produto
que permite coletar o máximo de aprendizado sobre
os usuários com o menor esforço possível. O objetivo
não é entregar o produto mais simples, mas sim
validar as hipóteses mais críticas do negócio antes
de investir no desenvolvimento completo.

No contexto da Medicano, o MVP contempla os fluxos
essenciais para validar a proposta de valor da
plataforma: cadastro e autenticação, triagem por
chatbot, busca por especialidade e localização,
agendamento de consultas e gestão de agenda.
Funcionalidades como notificações por e-mail,
pagamentos integrados, teleconsulta e a evolução
da arquitetura de IA para microserviços independentes
são documentadas como trabalhos futuros, a serem
implementadas em versões subsequentes da plataforma
\cite{ries2011}.
